\documentclass[11pt,a4paper]{article}
\usepackage[utf8]{inputenc}
\usepackage[T1]{fontenc}
\usepackage[french]{babel}
\usepackage{lmodern}
\usepackage{geometry}
\geometry{margin=2.3cm}
\usepackage{microtype}
\usepackage{setspace}
\setstretch{1.08}
\usepackage{xcolor}
\definecolor{cpblue}{HTML}{244A73}
\definecolor{cpgreen}{HTML}{2E7D6E}
\definecolor{cporange}{HTML}{C06B28}
\definecolor{cpgray}{HTML}{F3F5F7}
\definecolor{cpdark}{HTML}{24313B}
\usepackage{booktabs}
\usepackage{tabularx}
\usepackage{array}
\usepackage{longtable}
\usepackage{amsmath,amssymb}
\usepackage{enumitem}
\setlist{itemsep=2pt,topsep=4pt,parsep=0pt}
\usepackage{hyperref}
\hypersetup{colorlinks=true,linkcolor=cpblue,urlcolor=cpblue,citecolor=cpblue}
\usepackage{fancyhdr}
\pagestyle{fancy}
\fancyhf{}
\lhead{\small Note de transfert SCONTO-SVU}
\rhead{\small V1, \today}
\cfoot{\small\thepage}
\setlength{\headheight}{14pt}
\usepackage{titlesec}
\titleformat{\section}{\large\bfseries\color{cpblue}}{\thesection}{0.5em}{}
\titleformat{\subsection}{\normalsize\bfseries\color{cpdark}}{\thesubsection}{0.5em}{}
\titlespacing*{\section}{0pt}{14pt}{7pt}
\titlespacing*{\subsection}{0pt}{10pt}{5pt}
\newcommand{\code}[1]{\texttt{#1}}

\begin{document}

\begin{center}
{\Large\bfseries\color{cpblue} Note de transfert : modèles SCONTO-SVU\par}
\vspace{0.2cm}
{\normalsize Version générique, variante DataCo/Prophet/Recalibrator, fondation expérimentale\par}
\vspace{0.15cm}
{\small Note intermédiaire V1, document de travail\par}
\end{center}
\vspace{0.3cm}

\section{Objet du document}

Cette note décrit l'état des dépôts et modèles SCONTO-SVU au moment du
transfert, pour permettre à un développeur qui rejoint le projet de
comprendre rapidement quelles versions existent, laquelle est stable,
laquelle est spécialisée DataCo, laquelle est expérimentale, ce qui a été
corrigé, ce qui reste à faire, et ce qui peut ou non être fusionné à ce
stade.

\textbf{Cette note est une version intermédiaire (V1).} Les expérimentations
désignées EXP-0, E4 et E5 sont en cours ; aucune n'est présentée ici comme
définitivement validée. La branche expérimentale ne doit pas être fusionnée
dans le master stable avant les validations décrites section~\ref{sec:merge}.

\section{Cartographie des versions}
\label{sec:cartographie}

Deux dépôts Git et un ensemble de variantes locales du modèle DataCo sont
concernés.

\begin{center}
\small
\begin{tabularx}{\textwidth}{p{3.0cm}p{3.4cm}p{1.5cm}X}
\toprule
Dépôt & Branche & Commit & Rôle \\
\midrule
\code{sconto-svu-\allowbreak generic-\allowbreak anylogic-model} & \code{integration/\allowbreak claude-\allowbreak generic-merge} & \code{ba05f34} & Master générique stable, fusion figée \\
\code{sconto-svu-\allowbreak generic-\allowbreak anylogic-model} & \code{feature/\allowbreak experimentation-\allowbreak foundation} & \code{e94f5ec} & Fondation expérimentale (EXP-0), en validation \\
\code{anylogic-\allowbreak model-\allowbreak documentation} & \code{documentation-\allowbreak latex-\allowbreak rewrite} & \code{3698c14} & Documentation technique DataCo/Prophet/Recalibrator \\
\bottomrule
\end{tabularx}
\end{center}

Le modèle générique et la variante DataCo ne sont \textbf{pas} le même
artefact : le noyau générique ne contient ni Prophet, ni Recalibrator, ni
AutoCommande DataCo (section~\ref{sec:dataco-version}). Les variantes
DataCo identifiées dans l'espace de travail local sont cartographiées à la
section~\ref{sec:dataco-identite}.

\section{Version générique stable}

\textbf{Branche} : \code{integration/\allowbreak claude-generic-merge}. \textbf{Commit} :
\code{ba05f348\allowbreak 729ec607\allowbreak ae4716ab\allowbreak f97db1f7\allowbreak b4a7cce6}. \textbf{Artefact} :
\code{model/\allowbreak SCONTO\_SVU\_\allowbreak GENERIC\_MASTER.alp}.

\textbf{Statut} : fusion terminée, validée par build et runs utilisateur,
figée avant toute expérimentation. Ce commit est le point de départ exact
de la branche expérimentale décrite section~\ref{sec:exp0} ; aucun
changement n'y a été apporté depuis.

Les blocs suivants ont été intégrés et validés (build AnyLogic et runs
utilisateur) au cours de la fusion :

\begin{itemize}
  \item architecture générique pilotée par scénarios JSON (aucune
        configuration métier codée en dur dans le moteur) ;
  \item hiérarchie SCOR (processus Source, Make, Deliver, Return, niveaux
        P1 à P5) ;
  \item cartographie VSM des micro-activités et des flux ;
  \item architecture multi-agents holonique (ADACOR/PROSA) : agents
        stratégique, tactiques, coordinateurs, opérationnels ;
  \item hiérarchie ISA-95 et messagerie AER (tableau noir partagé) ;
  \item pipeline de Performance Index (méthode Theeranuphattana/Chan-Qi
        adaptée, agrégation floue) et indicateurs KPI par poste ;
  \item politiques de stock MTS, MTO et ETO ;
  \item gestion multi-commandes et multi-produit ;
  \item réapprovisionnement autonome par produit (ordres \code{REAPPRO\_*}) ;
  \item ordonnancement réel par insertion FIFO dans \code{actionsMetierDifferees}
        (section~\ref{sec:corrections}) ;
  \item exports Excel et ontologiques (ABox RDF/Turtle).
\end{itemize}

E4 et E5 (section~\ref{sec:e4e5}) ne sont \textbf{pas} implémentées dans
cette version : aucune variable de décision, aucun espace de configurations,
aucune boucle Pareto/AHP de campagne d'optimisation n'existe dans ce
commit.

\section{Version DataCo / Prophet / Recalibrator}
\label{sec:dataco-version}

La variante DataCo est une extension spécialisée du modèle générique,
maintenue comme un artefact distinct. Elle contient notamment :

\begin{itemize}
  \item Prophet, prévision mensuelle de référence entraînée sur la période 2015 à 2016 ;
  \item l'AutoCommande DataCo, qui rejoue la demande 2017 sous forme de
        commandes clientes journalières ;
  \item le Recalibrator, qui combine Prophet et la demande observée via un
        facteur de confiance $\alpha_j = \sqrt{j/N}$ (confirmé dans le code :
        \code{double alpha = Math.sqrt(Math.max(0.0, Math.min(1.0,
        (double) jour / (double) dernierJour)));}) ;
  \item le fonctionnement multi-produit DataCo (familles \code{DATACO\_\allowbreak SPORTS},
        \code{DATACO\_\allowbreak CLOTHING}, \code{DATACO\_\allowbreak ELECTRONICS}, configurées par
        scénario JSON, non codées en dur dans le moteur générique) ;
  \item les prévisions mensuelles et la recalibration hebdomadaire ;
  \item les exports de résultats (Excel, ABox) associés aux campagnes
        DataCo.
\end{itemize}

\textbf{Le modèle générique ne contient pas Prophet ni Recalibrator.} Cette
séparation est volontaire et documentée dans les instructions de fusion du
dépôt : le noyau générique peut recevoir des capacités génériques utiles à
DataCo (multi-produit générique, stocks par produit, interfaces
d'extension), mais aucune connaissance métier DataCo codée en dur. Les
fonctions \code{genererCommandeDataCoPourDate()},
\code{recalibrerPrevisionPourDate()} et les structures
\code{prophetForecastsOriginal} sont absentes du master générique
(\code{ba05f34}) ; elles n'existent que dans les variantes DataCo
répertoriées section~\ref{sec:dataco-identite}.

\textbf{Ne pas injecter Prophet/Recalibrator dans le master générique
pendant une opération de fusion sans décision architecturale explicite.}

\subsection{Résultats DataCo de référence}

Les valeurs suivantes proviennent du run de référence RUN~1
(\code{simToRealSeconds = 3600}), documenté dans
\code{documentation/assets/dataco\_final\_run/RUN\_REFERENCE.md} du dépôt
de documentation.

\begin{center}
\small
\begin{tabular}{lr}
\toprule
Indicateur & Valeur \\
\midrule
Commandes générées & 365 \\
Demande totale & 99\,896 \\
Commandes closes & 79 \\
dont closes en retard & 20 \\
Commandes EN\_ATTENTE & 286 \\
Quantité livrée au snapshot & 22\,469 \\
Performance Index au snapshot & 8,871/10 \\
\bottomrule
\end{tabular}
\end{center}

\begin{center}
\small
\begin{tabular}{lrr}
\toprule
Métrique & Prophet & Recalibrator \\
\midrule
MAE & 2\,491,89 & 123,06 \\
MAPE & 45,75\,\% & 1,91\,\% \\
RMSE & 3\,441,15 & 180,71 \\
\bottomrule
\end{tabular}
\end{center}

Réduction de l'erreur absolue totale entre Prophet et Recalibrator :
environ 95,06\,\%.

\textbf{RUN~1 n'est pas terminal.} 286 commandes restent EN\_ATTENTE au
snapshot ; le Performance Index de 8,871/10 est un indicateur intermédiaire,
pas une performance logistique terminale. Ces valeurs ne doivent pas être
présentées comme un résultat final de la campagne.

\section{Identification de la variante DataCo réellement exécutée}
\label{sec:dataco-identite}

Onze fichiers \code{.alp} liés à DataCo/Prophet/Recalibrator ont été
localisés dans l'espace de travail. Pour chacun : chemin, taille, SHA-256,
présence des fonctions \code{genererCommandeDataCoPourDate()} et
\code{recalibrerPrevisionPourDate()}, et statut.

\begin{center}
\small
\begin{tabularx}{\textwidth}{p{4.0cm}rp{1.3cm}X}
\toprule
Fichier & Taille (o) & Fonctions DataCo & Statut \\
\midrule
\code{sources/\allowbreak model/\allowbreak SCONTO\_\allowbreak SVU\_FINAL\_\allowbreak VSM\_FIX\_\allowbreak CANDIDATE.alp} & 2\,237\,244 & Absentes & Candidat VSM antérieur à l'intégration DataCo ; ne contient ni Prophet ni AutoCommande DataCo \\
\code{...\_DATACO\_\allowbreak COMPLETE.alp} (migration) & 2\,338\,928 & Présentes & Étape de développement DataCo (2026-09-09) \\
\code{...\_DATACO\_\allowbreak AUTOMATED.alp} (migration) & 2\,336\,084 & Présentes & Étape de développement (2026-09-09) \\
\code{...\_DATACO\_\allowbreak AUTOMATED\_\allowbreak CLOCK.alp} (migration) & 2\,338\,176 & Présentes & Étape de développement (2026-09-09) \\
\code{...\_DATACO\_\allowbreak DYNAMIC\_\allowbreak GRAPHS.alp} (migration) & 2\,334\,561 & Présentes & Étape de développement (2026-09-11) \\
\code{...\_DATACO\_\allowbreak YEAR\_\allowbreak BOUNDED\_\allowbreak FINAL.alp} (migration) & 2\,237\,392 & Présentes & Étape de développement (2026-09-11) \\
\code{...\_DATACO\_\allowbreak BARCHARTS\_\allowbreak REALISTIC\_\allowbreak RECAL.alp} (migration) & 2\,367\,526 & Présentes & Étape de développement (2026-09-11) \\
\code{...\_DATACO\_\allowbreak MULTIPRODUCT\_\allowbreak FOUNDATION.alp} (migration) & 2\,376\,717 & Présentes & Étape de développement (2026-09-12) \\
\code{...\_DATACO\_\allowbreak MULTIPRODUCT\_\allowbreak CONCURRENT\_\allowbreak FIX.alp} (migration, 2026-09-15) & 2\,371\,853 & Présentes & \textbf{Identique octet pour octet (hors nom interne) à \code{SCONTO\_SVU\_\allowbreak DATACO\_FINAL.alp}, voir ci-dessous} \\
\code{reference/\allowbreak dataco-\allowbreak multiproduct/\allowbreak ...\_DATACO\_\allowbreak MULTIPRODUCT\_\allowbreak CONCURRENT\_\allowbreak FIX.alp} (dépôt générique, 2026-09-16) & 2\,244\,316 & Présentes & Même nom que la ligne précédente, contenu substantiellement différent (voir remarque) \\
\code{SCONTO\_SVU\_\allowbreak LIVRAISON\_\allowbreak ENCADREMENT/\allowbreak 01\_MODELES/\allowbreak SCONTO\_SVU\_\allowbreak DATACO\_FINAL.alp} & 2\,371\,806 & Présentes & \textbf{Fichier co-localisé avec les exports RUN~1/RUN~2 ; retenu comme modèle de référence de la campagne DataCo documentée} \\
\bottomrule
\end{tabularx}
\end{center}

\subsection{Résultat de la comparaison}

La comparaison octet par octet entre \code{SCONTO\_SVU\_\allowbreak DATACO\_FINAL.alp}
et la variante \code{...\_DATACO\_\allowbreak MULTIPRODUCT\_\allowbreak CONCURRENT\_FIX.alp} du
dossier de migration (datée du 2026-09-15) montre une différence unique :
le champ interne \code{<Name>} du projet AnyLogic
(\code{SCONTO\_SVU\_\allowbreak DATACO\_FINAL} contre
\code{SCONTO\_SVU\_\allowbreak FINAL\_\allowbreak VALIDATED\_\allowbreak FORECAST\_\allowbreak DATACO\_\allowbreak MULTIPRODUCT\_\allowbreak CONCURRENT\_FIX}).
Tout le reste du fichier est identique. \textbf{Conclusion démontrée, non
supposée} : \code{SCONTO\_SVU\_\allowbreak DATACO\_FINAL.alp} est une copie renommée de
cette variante, sans aucune autre modification.

Le fichier de même nom présent dans
\code{reference/dataco-multiproduct/} du dépôt générique, daté du
lendemain (2026-09-16), est en revanche substantiellement différent
(taille inférieure d'environ 127\,Ko, différences réparties sur l'ensemble
du fichier) : il s'agit d'un artefact distinct malgré le nom identique.
Cette homonymie est une source de confusion à corriger (renommage ou
suppression d'une des deux copies) avant toute reprise du projet.

\textbf{Le champ runtime \code{modeleCandidate}}, présent dans le manifeste
exporté par les campagnes DataCo, est une étiquette codée en dur dans le
générateur de manifeste (confirmée identique dans plusieurs variantes du
dépôt) et \textbf{ne constitue pas une preuve} du fichier réellement
exécuté ; il ne doit pas être utilisé comme tel. L'identification
ci-dessus repose exclusivement sur la comparaison de contenu et sur la
co-localisation du fichier avec les exports RUN~1/RUN~2.

\section{Principales évolutions du modèle}

\subsection{Ligne générique}

Fusion progressive de \code{model/SCONTO\_SVU\_GENERIC\_MASTER.alp} avec
les apports du modèle collègue, en six blocs fonctionnels (A, M, B, C, D,
E), chacun audité, corrigé, validé statiquement puis par build et run
utilisateur avant de passer au suivant. Clôture du dernier bloc (E,
ordonnancement) et gel du master au commit \code{ba05f34}.

\subsection{Ligne DataCo}

Progression documentée par une suite de fichiers \code{.alp} successifs
(2026-09-09 à 2026-09-15) : intégration initiale, automatisation de
l'horloge, graphiques dynamiques, bornage calendaire de l'année cible,
recalibration réaliste, fondation multi-produit, puis correction de
concurrence multi-produit, ce dernier état étant celui effectivement
exécuté pour produire RUN~1/RUN~2 (section~\ref{sec:dataco-identite}).

\subsection{Fondation expérimentale}

Trois passes successives sur la branche \code{feature/experimentation-foundation},
détaillées section~\ref{sec:exp0}.

\section{Corrections fonctionnelles importantes}
\label{sec:corrections}

\begin{center}
\small
\begin{tabularx}{\textwidth}{p{3.4cm}Xp{3.6cm}}
\toprule
Problème & Correction & Effet \\
\midrule
Commandes MTO/ETO client bloquées en boucle infinie \code{ANALYSE\_STOCK} (garde de sécurité qui ne couvrait pas ces types) & Séparation stricte MTS / MTO-ETO dans \code{analyserStockCommandeOrchestree()} : après les deux branches MTS inchangées, appel direct à \code{analyserMatiereCommandeOrchestree()}, déjà utilisée sans modification par le réapprovisionnement autonome & Les commandes MTO/ETO progressent jusqu'à clôture ; garde de réentrance étendue pour éviter une double production en cas de réveil concurrent \\
Ordonnanceur dédié (\code{ordonnancer\allowbreak Production()}, tri priorité/EDD) jamais exécuté utilement & Audité et confirmé inactif : \code{carnetCommandesMTO} n'est jamais peuplé dans tout le master (0 appel \code{.add()}) ; clarifié comme mécanisme hérité non actif, non réactivé & La politique réelle d'ordonnancement, FIFO par insertion dans \code{actionsMetierDifferees}, est documentée comme telle plutôt que supposée \\
Association flux/commande imprécise (\code{CMD\_4} pouvait matcher \code{CMD\_44}) & Matching strict de préfixe & Aucune confusion entre commandes aux identifiants proches \\
Engagement client (délai promis \code{tPromise}) écrit deux fois, dont un calcul mort avant résolution complète du scénario & Écriture prématurée supprimée ; seule écriture conservée, après résolution du scénario et du type de commande ; delai par défaut et par produit rendus configurables (\code{champDelaiPromesseParDefautSec}, \code{ScenarioFlux.promisedDelaySec}) & Une seule source de vérité pour \code{tPromise}, configurable sans modification du moteur \\
Détection des retards clients non centralisée & \code{commandeEstTerminale()} établi comme prédicat canonique unique (\code{SERVIE}/\code{EN\_RETARD}/\code{REFUSEE}/\code{BLOQUE\_*}) ; \code{retardConstate} et \code{modeEngagementClient} comme champs dédiés, détection continue via événement périodique & Une seule définition de la terminalité et du retard, réutilisée par tous les compteurs (KPI, exports, diagnostics) \\
Compteur de commandes ouvertes potentiellement ambigu (client vs réapprovisionnement) & Vérifié par lecture directe : exclut déjà les ordres \code{REAPPRO\_*} via \code{estOrdreStockAutonome()} & Compteur \code{nombreCommandesOuvertes()} correct pour un usage \og commandes clientes uniquement\fg{} sans modification nécessaire \\
Incohérences RL/RS du pipeline SCOR vers PI (proxies non calibrés, métrique retournant 0 malgré stock non nul, warm-up affichant PI=0) & Corrections ciblées bloc par bloc (B.1 à B.3-FIX2), sources et unités alignées, proxies explicitement documentés comme tels & Pipeline PI/SCOR revalidé par build et runs utilisateur, clôturé et figé \\
\bottomrule
\end{tabularx}
\end{center}

\section{Architecture SCOR / VSM / SMA conservée}

L'architecture retenue à l'issue de la fusion reste inchangée par les
travaux EXP-0 :

\begin{itemize}
  \item \textbf{SCOR} : processus Source, Make, Deliver, Return, codes de
        niveau 1 à 3, catalogue de métriques associé à chaque code ;
  \item \textbf{VSM} : micro-activités et flux physiques/informationnels
        cartographiés par scénario JSON ;
  \item \textbf{SMA holonique (ADACOR/PROSA)} : un agent stratégique
        (\code{StrategicAgent}), quatre agents tactiques (un par
        macro-processus SCOR principal, \code{TacticalAgent}), des
        coordinateurs dédiés par macro-processus (\code{CoordinatorAgent})
        et des agents opérationnels par poste (\code{OperationalAgent}) ;
  \item \textbf{ISA-95} : hiérarchie industrielle conservée dans l'export
        ontologique ;
  \item \textbf{AER} : messagerie d'alerte/escalade via un tableau noir
        partagé (\code{BlackboardAgent}), remontée opérationnel $\to$
        tactique $\to$ supervision ;
  \item \textbf{Performance Index} : méthode Theeranuphattana/Chan-Qi
        adaptée, agrégation floue des cinq attributs (RL, RS, AG, CO, AM).
\end{itemize}

\section{Expérimentations et instrumentation}
\label{sec:exp0}

Branche \code{feature/experimentation-foundation}, créée à partir de
\code{ba05f34}. Trois commits, chacun documenté dans
\code{EXPERIMENTATION\_FOUNDATION.md} :

\begin{itemize}
  \item \textbf{\code{56f763f}}, fondation initiale : audit exhaustif des
        sources d'aléatoire (toutes passent par
        \code{getDefault\allowbreak RandomGenerator()}) ; paramètre
        \code{seedExperiment} appliqué en tête de \code{demarrer\allowbreak Simulation()} ;
        distinction entre fin de génération de la demande et état terminal
        opérationnel ; diagnostic terminal structuré ; champs de
        traçabilité additifs dans l'export ontologique.
  \item \textbf{\code{82f60d2}}, vérification de trois ambiguïtés :
        confirmation que le compteur de commandes ouvertes exclut déjà les
        ordres de réapprovisionnement ; confirmation que la définition de
        fin de génération couvre le scénario de choc de demande utilisé
        comme référence pour une future campagne d'optimisation ;
        correction de la persistance de la graine dans le scénario JSON
        (sérialisation en chaîne plutôt qu'en nombre à virgule flottante,
        pour éviter une perte de précision).
  \item \textbf{\code{e94f5ec}}, correction d'un défaut de
        reproductibilité réel, mis en évidence par un test runtime à
        graine identique : plusieurs événements métier cycliques étaient
        phasés sur le temps absolu du moteur AnyLogic plutôt que sur le
        début logique du run. Huit événements rephasés via l'API standard
        \code{restart()}, sans modification d'aucune règle métier, période
        ou paramètre de panne.
\end{itemize}

\textbf{Statut : en cours de validation runtime. Ne pas fusionner dans le
master stable à ce stade.}

\section{Reproductibilité}

Le test runtime effectué après le commit \code{56f763f} a montré que,
même avec \code{seedExperiment=1001} identique sur deux runs, les
trajectoires divergeaient à partir du traitement d'un ordre de
réapprovisionnement autonome. La cause exacte, confirmée par lecture du
code, est un événement cyclique (\code{cycleArbitrage}, agent tactique)
configuré pour se déclencher sur une grille de temps absolue du moteur,
indépendamment de l'instant où l'utilisateur démarre la campagne.

La correction appliquée au commit \code{e94f5ec} rephase huit événements
métier identifiés comme pertinents (mutation d'état ou décision affectant
le flux) relativement au début logique du run.

\textbf{Point de vigilance explicite} : cette correction n'a pas encore été
validée par un nouveau test runtime à deux graines identiques après
\code{e94f5ec}. Statut exact : \textbf{implémentée, build à valider,
runtime à revalider.}

\section{Travaux encore en cours}

\begin{itemize}
  \item Revalidation runtime de la correction de rephasage
        (section précédente) : deux runs à \code{seedExperiment} identique
        doivent produire des trajectoires identiques après normalisation
        au début du run.
  \item Un état latent distinct a été identifié sans être corrigé :
        \code{tProchainePanne} et \code{enPanne} (par poste) ne sont pas
        explicitement réinitialisés entre plusieurs runs successifs dans
        une même session AnyLogic. Non observé dans le test effectué,
        probablement parce que la valeur initiale restait à son défaut
        avant le premier run ; à surveiller si une session enchaîne
        plusieurs runs sans rouvrir le modèle.
  \item Le problème \code{demandGenerationFinished} observé à la clôture
        finale du run (\code{RUN\_FINALIZED}) a été signalé par
        l'utilisateur mais n'a pas encore été investigué : aucune
        dépendance technique avec la correction de rephasage n'a été
        démontrée, traitement reporté à une passe distincte.
  \item E4 et E5 (section~\ref{sec:e4e5}) ne sont pas exécutées.
\end{itemize}

\section{E4 et E5, travaux futurs}
\label{sec:e4e5}

\textbf{Aucune de ces campagnes n'est exécutée à ce jour.} Elles sont
présentées ici uniquement comme protocoles formalisés dans la
documentation, à titre de travaux futurs s'appuyant sur la fondation
expérimentale EXP-0.

\textbf{E4, optimisation ZENER.} Variables de décision : $x_1$ (capacité
du poste critique), $x_2$ (taille des sous-lots), $x_3$ (règle de
priorité), $x_4$ (politique de stock), soit 144 configurations candidates.
Méthode : exploration factorielle, front de Pareto, puis sélection AHP
parmi les solutions Pareto-efficaces.

\textbf{E5, optimisation DataCo fondée sur la prévision.} Compare quatre
politiques : P0 (réactive, sans prévision), P1 (Prophet seul), P2 (Prophet
et Recalibrator), P3 (Prophet, Recalibrator et paramètres de planification
optimisés). Variables de décision par famille simulée :
$x_p = (SS_p, s_p, Q_p, f_p)$.

\section{État actuel des branches Git}

\begin{center}
\small
\begin{tabularx}{\textwidth}{p{4.6cm}p{4.2cm}X}
\toprule
Dépôt / branche & HEAD & Statut \\
\midrule
\code{sconto-svu-generic-\allowbreak anylogic-model} \newline \code{integration/\allowbreak claude-generic-merge} & \code{ba05f34} & Stable, figée, fusion terminée \\
\code{sconto-svu-generic-\allowbreak anylogic-model} \newline \code{feature/\allowbreak experimentation-foundation} & \code{e94f5ec} & Expérimentale, en validation runtime, non fusionnée \\
\code{anylogic-model-\allowbreak documentation} \newline \code{documentation-\allowbreak latex-rewrite} & \code{3698c14} & Documentation DataCo/Prophet/Recalibrator à jour \\
\bottomrule
\end{tabularx}
\end{center}

\section{Instructions de merge pour le collègue}
\label{sec:merge}

\textbf{A. Pour récupérer le modèle générique stable} : \code{checkout}
\code{integration/\allowbreak claude-\allowbreak generic-merge} au commit \code{ba05f34}.
Artefact : \code{model/\allowbreak SCONTO\_SVU\_\allowbreak GENERIC\_\allowbreak MASTER.alp}.

\textbf{B. Ne pas fusionner} \code{feature/experimentation-foundation}
pour l'instant : le runtime est encore en cours de validation
(section~\ref{sec:exp0}).

\textbf{C. Conserver la variante DataCo séparée} du master générique.

\textbf{D. Ne pas injecter} Prophet ni Recalibrator dans le master
générique pendant une opération de fusion sans décision architecturale
explicite.

\textbf{E. Avant tout merge futur de la branche expérimentale}, valider
dans l'ordre :

\begin{enumerate}
  \item build AnyLogic 8.9.5 sans erreur ;
  \item reproductibilité à \code{seedExperiment=1001} (deux runs, mêmes
        trajectoires après normalisation) ;
  \item divergence attendue à \code{seedExperiment=1002} ;
  \item les trois cas de terminalité (génération en cours, génération
        terminée avec encours, terminal atteint) ;
  \item réinitialisation de l'état des pannes entre runs successifs ;
  \item non-régression des scénarios MTS, MTO, ETO, multi-commandes,
        multi-produit déjà validés sur le master.
\end{enumerate}

\section{Points de vigilance}

\begin{itemize}
  \item Deux fichiers portent le même nom (\code{...\_DATACO\_MULTIPRODUCT\_CONCURRENT\_FIX.alp})
        dans deux emplacements différents, avec un contenu substantiellement
        différent (section~\ref{sec:dataco-identite}) : à clarifier avant
        toute reprise.
  \item Le champ runtime \code{modeleCandidate} ne doit jamais être utilisé
        comme preuve d'identité d'un fichier exécuté.
  \item Les exports runtime volumineux (ABox TTL jusqu'à plusieurs Mo) ne
        sont pas committés dans le dépôt de documentation ; ils restent
        disponibles dans le dossier de livraison local et sont nommés dans
        \code{RUN\_REFERENCE.md}.
  \item Le champ exporté \code{entitiesInProcessAtPosts} (EXP-0) compte des
        occupations dans les structures de traitement des postes, pas
        garanti comme un nombre d'agents métier uniques.
  \item Le modèle utilise un seul flux aléatoire partagé pour tous les
        mécanismes stochastiques : deux configurations qui consomment un
        nombre différent de tirages avant un point donné désynchronisent le
        flux au-delà de ce point, même à graine identique. Une séparation
        en générateurs dédiés par mécanisme resterait à faire pour un
        appariement strict des nombres aléatoires communs, si les
        campagnes E4/E5 le requièrent.
\end{itemize}

\section{Tableau de traçabilité des artefacts}

\begin{center}
\small
\begin{tabularx}{\textwidth}{p{3.2cm}p{2.6cm}p{1.6cm}Xcp{1.6cm}}
\toprule
Artefact & Branche/dépôt & Commit & Rôle & Merge & Remarques \\
\midrule
\code{SCONTO\_SVU\_\allowbreak GENERIC\_\allowbreak MASTER.alp} & integration/claude-\allowbreak generic-merge & ba05f34 & Master générique & Oui & Stable, figé \\
\code{SCONTO\_SVU\_\allowbreak DATACO\_FINAL.alp} & hors dépôt structuré (livraison locale) & n/d & Modèle DataCo exécuté pour RUN~1/RUN~2 & Non applicable & Identique (hors nom) à la variante migration 2026-09-15 \\
\code{...\_VSM\_FIX\_\allowbreak CANDIDATE.alp} & hors dépôt structuré & n/d & Candidat antérieur, sans DataCo & Non & Historique \\
\code{feature/\allowbreak experimentation-\allowbreak foundation} (branche) & sconto-svu-\allowbreak generic & e94f5ec & Fondation EXP-0 & Non, en validation & Revalidation runtime requise \\
\code{Documentation\_\allowbreak Module\_\allowbreak Prevision\_\allowbreak SCONTO\_SVU\_\allowbreak corrige.tex} & anylogic-model-\allowbreak documentation & 3698c14 & Documentation module DataCo & n/a & À jour \\
\code{REPONSES\_\allowbreak ENCADREMENT\_\allowbreak PROPHET\_ALPHA\_\allowbreak OPTIMISATION.tex} & anylogic-model-\allowbreak documentation & 3698c14 & Synthèse académique Prophet/\allowbreak$\alpha_j$/\allowbreak optimisation & n/a & À jour \\
\code{RUN\_\allowbreak REFERENCE.md} & anylogic-model-\allowbreak documentation & 3698c14 & Traçabilité RUN~1/RUN~2 & n/a & À jour \\
\bottomrule
\end{tabularx}
\end{center}

\section{Recommandations pour la reprise par le collègue}

\begin{enumerate}
  \item Partir du master générique (\code{ba05f34}) pour toute nouvelle
        fonctionnalité générique ; ne pas repartir de la branche
        expérimentale tant qu'elle n'est pas revalidée.
  \item Ne pas modifier \code{model/SCONTO\_SVU\_GENERIC\_MASTER.alp}
        directement sans créer une branche dédiée, conformément au
        workflow déjà établi dans le dépôt.
  \item Pour toute reprise du travail DataCo, utiliser
        \code{SCONTO\_SVU\_DATACO\_FINAL.alp} comme référence, en gardant à
        l'esprit qu'il s'agit d'une copie renommée de la variante de
        migration datée du 2026-09-15, et non d'un artefact distinct.
  \item Avant de fusionner \code{feature/experimentation-foundation},
        exécuter le protocole de validation de la section~\ref{sec:merge}.E
        dans son intégralité.
  \item Clarifier l'homonymie signalée section~\ref{sec:dataco-identite}
        avant toute automatisation qui se fierait au nom de fichier plutôt
        qu'à son contenu.
  \item Ne pas exécuter E4 ou E5 avant que la reproductibilité et la
        condition terminale de EXP-0 soient entièrement validées.
\end{enumerate}

\section*{Annexe : chronologie des principales évolutions}
\addcontentsline{toc}{section}{Annexe : chronologie des principales évolutions}

\begin{itemize}
  \item Fusion générique, blocs A (corrections sûres) et M (fondation
        multi-produit).
  \item Bloc B : refonte du pipeline Performance Index / SCOR.
  \item Bloc C : détection continue des retards, engagement client
        configurable.
  \item Bloc D : audit du moteur d'exécution et des exports.
  \item Bloc E : audit et correction de l'ordonnancement, raccordement du
        chemin de production MTO/ETO.
  \item Clôture et gel du master générique (\code{ba05f34}).
  \item Développement DataCo : Prophet, AutoCommande, Recalibrator,
        multi-produit DataCo (progression de fichiers \code{.alp}
        successifs, 2026-09-09 à 2026-09-15).
  \item Documentation technique DataCo/Prophet/Recalibrator (dépôt de
        documentation, plusieurs passes de rigueur scientifique).
  \item EXP-0 : graine expérimentale, condition terminale, diagnostic.
  \item EXP-0.1 : vérification de trois ambiguïtés (commandes ouvertes,
        couverture E4, précision du round-trip JSON de la graine).
  \item EXP-0.2 : correction de la reproductibilité runtime (rephasage des
        événements métier cycliques).
\end{itemize}

\end{document}
